\section{Metodología de desarrollo}

\subsection{Selección de una metodología ágil}

Se a decidido que para el desarrollo de este proyecto se utilizará la metodología Scrum, esta metodología altamente conocida por su enfoque iterativo incremental promueve la adaptabilidad y la entrega continua, de esto se desligan distintas razónes de porque fue seleccionada:

\begin{enumerate}
    
    \item \textbf{Iteraciones cortas (Sprints):} Los ciclos cortos de Scrum son la clave para este proyecto, el poder contar con versiones funcionales desde etapas tempranas permitirá medir el avance y el valor que esta entregando la solución casi desde el principio, y así evaluar la viabilidad de posibles incrementos.

    \item \textbf{Retroalimentación temprana:} Siguiendo con lo del punto anterior, estas versiones funcionales casi al inicio del ciclo de vida del proyecto permitirá realizar pruebas que permitan detectar de forma oportuna errores o desviaciones, mejorando la calidad del producto final.
    
    \item \textbf{Alta adaptabilidad a cambios:} La posibilidad mencionada anteriormente de recibir retroalimentación continua sobre el producto permitirá adaptarse a cualquier cambio que sea necesario que fuera imprevisto durante la planificación inicial. Esto es clave teniendo en cuenta que es el primer proyecto basado en LLM de los integrantes del grupo, por lo que se podrían pasar cosas por alto en una planificación inicial por mera falta de experiencia en el área.
\end{enumerate}


\subsection{Planificación de iteraciones y definición de roles en el equipo}

\subsubsection{Plan de iteraciones}

a

\subsubsection{Roles del equipo y tareas asociadas}

\begin{enumerate}
    \item Nicolás Sánchez - Scrum Master: Encargado de facilitar la adopción de la metodología Scrum por parte del equipo y asegurando el cumplimiento de las prácticas y entregas asociadas a dicha metodología, además de coordinar el flujo de trabajo y reuniones grupales. 
    
    \item Matías Hernández - Product Owner: Encargado de definir los requerimientos del producto, validar el cumplimiento de este hacía las necesidades de los posibles usuarios desde la toma de decisiones sobre los pasos a seguir para obtener un buen producto, hasta que este sea desplegado para los usuarios finales.
    
    \item Desarrollador backend - Sergio Espinoza: Encargado del diseño e implementación del backend, base de datos e integración de módulos.
    
    \item Desarrollador frontend - Daniel Eguiluz: Encargado del diseño e implementación de la interfaz para los usuarios y para el consumo del backend.
\end{enumerate}

Para la selección de estos roles y respectivas tareas se tomaron en cuenta los trabajos previos realizados por cada miembro, de forma que vaya acorde a sus áreas de especialización y así asegurar el cumplimiento de los objetivos del proyecto.


\subsection{Diagrama de usos éticos de la IA (MyAI Compass)}

a

\subsection{Repositorio de control de versiones}

Para el control de versiones del código y los informes del proyecto se hará uso de github, específicamente del repositorio presente en el siguiente enlace: \url{https://github.com/Ch3chS/TAVI-Charlie}.